\chapter{10. 安全与应急管理}\label{ux5b89ux5168ux4e0eux5e94ux6025ux7ba1ux7406}

安全体系是``日轮之冕''得以长期稳定运行的保障。我们搭建安全与应急管理体系的核心，是尽力去追求 ``零风险'' 的绝对完美，即便不能做到，也会让所有风险始终处于可感知、可控制的范围之内。在复杂未知的环境中，这套体系更看重社会系统的持续运转与快速恢复能力 ------AI与智能机器人是我们重要的辅助伙伴，负责监测、分析和提供建议与代替人工劳动力。尽管AI是我们世界的重要组成部分，但关键处置权永远握在制度与人工手中，以避免技术失效或误判带来更大风险。\\
我们坚持 ``必要够用'' 原则，不盲目堆砌算力和设备：日常低能耗运行，仅在需要时集中调用资源；不搞无差别管控，既保障安全底线，也留足个人自由与生活温度。针对无人愿意自愿戍边、当兵的情况，我们构建了 ``机器人为主、制度兜底'' 的防御模式，让安全守护既高效又贴合大家的意愿。

\section{10.1 个人安全}\label{ux4e2aux4ebaux5b89ux5168}

个人安全的核心是 ``防患于未然''，我们不想做过度干预生活的 ``管家''，而是默默守护的 ``后盾''。

\subsection{10.1.1 日常防护赋能}\label{ux65e5ux5e38ux9632ux62a4ux8d4bux80fd}

居民都有专属 AI 代表，它和 ``长征'' 计算中枢实时联动，7×24 小时帮你评估身边的安全风险，哪怕是在 VR/AR 创作或游憩时也不例外。我们为大家配备了轻便的生命体征监测模块，能实时追踪心率、血氧和体温，一旦出现异常，会第一时间同步给 ``长征'' 中枢，让危险早被发现。而在大家常去的 Workshop 创作空间，建德院还安排了智能安防机器人，它们会主动识别火源、电气隐患这些问题，发现后立刻同步给启真院处理，让我们在创作和活动时更安心，下面是所需要使用的相关设备与技术：

1.可穿戴生命体征监测模块（参考：Huawei TruSeen 5.5+）\\
功能：实现对人体心率、血氧、体温的实时监测，同时危险时自动同步到``长征''中枢

2.智能安防机器人（参考：Boston Dynamics Spot）\\
功能：可以自动识别火源、电气隐患、人员跌倒等风险

3.VR/AR 风险识别系统（参考：Meta Quest Pro）\\
功能：在居民 MR 创作 / Workshop 内提前识别可能伤害轨迹，为居民提供更好的安全保障\\
来源：1.https://www.huawei.com2.https://bostondynamics.com/products/spot 3.https://www.meta.com/quest/

\subsection{10.1.2 突发风险应对}\label{ux7a81ux53d1ux98ceux9669ux5e94ux5bf9}

一旦遇到紧急情况，个人 AI 触发求救后，玉湖院部署的医疗辅助机器人 1 分钟内就能赶到身边，做急救、固定伤处、自动给氧这些基础处理，同时会调度运输机器人送医。如果是在海上遇险，AI 代表会联动海岸安防系统，救援无人机和机器人编队 24 小时待命，再恶劣的天气也能快速响应。下面是所需要的一些设备：

1.医疗辅助机器人（参考：F\&P Robotics P-Rob）\\
功能：可以做到医学急救、固定伤处、自动给氧等基础急救需求

2.海岸安防雷达（IALA X-band）\\
作用：24小时监控小平岛海面，以查看有无突发情况

来源： 1.https://www.fp-robotics.com/robots/p-rob/\\
2.https://iala-aism.org

\subsection{10.1.3 安全技能普及}\label{ux5b89ux5168ux6280ux80fdux666eux53ca}

安全技能不应被遗弃，而是应当成为``日轮之冕''居民的必备能力。建德院通过 MR 游憩场景开设安全演练，大家可以在自由时间里，模拟火灾疏散、核泄漏避难、落水逃生这些场景，为儿童提供娱乐的学习氛围。启真院还会按计划给每户分配智能应急包，机器人管家每月会检查物资是否齐全，确保遇到情况时能派上用场。所采用的MR 安全演练系统（参考 Varjo XR-4）https://varjo.com

\section{10.2 建筑与生活场所安全}\label{ux5efaux7b51ux4e0eux751fux6d3bux573aux6240ux5b89ux5168}

我们的建筑和生活场所，靠的是 ``长期监测'' 而非 ``频繁维修''，哪怕城市不断扩建，安全管理也不会失控。传感器长期、低能耗运行，用来监测建筑的变形、裂缝和环境异常，一旦出现问题就提前预警，而工程机器人主要在施工和检修时工作，平时不占用大量能源。这样可以保证城市扩建后，建筑数量增加，但安全管理不会失控。

\subsection{10.2.1 建筑智能适配}\label{ux5efaux7b51ux667aux80fdux9002ux914d}

遵义院主导建设的所有建筑，都按 ``防台风 + 抗地震 + 核安全'' 的高标准打造。建筑里内置了和 ``长征'' 中枢联动的监测传感器，能实时检测结构应变、裂缝、震感这些关键数据，一旦出现异常就提前预警。施工时还有 AI 视觉质检系统把关，质量合格后会报备人民大会监督，住得放心才是根本。

技术设备：1.结构健康监测传感器（参考：HBM 光纤 FBG 传感器）2.AI 视觉质检系统（NVIDIA Jetson Orin）\\
来源：1.https://www.hbm.com2.https://www.nvidia.com

\subsection{10.2.2 建筑消防安全管理}\label{ux5efaux7b51ux6d88ux9632ux5b89ux5168ux7ba1ux7406}

启真院统筹的智能消防机器人 24 小时在公共区域巡逻，发现火源能自动识别并喷洒灭火剂。公共区域的消防设施都接入了物联网，任何异常都能被及时发现。建德院也会通过公共媒体推送安全知识，让全民参与到消防安全保障中。

相关设备：1.消防巡逻机器人（参考：Unitree B1 搭载热像系统）2.公共区域物联网烟雾/气体监测器（Honeywell）\\
来源：1.https://www.unitree.com2.https://www.honeywell.com

\subsection{10.2.3 日常维护机制}\label{ux65e5ux5e38ux7ef4ux62a4ux673aux5236}

遵义院的维修机器人每天都会巡检建筑和核能管线，发现隐患后，数据会同步到 ``长征'' 中枢生成修复方案。如果家里或工作场所需要维修，我们通过 AI 提交申请，启真院 1 小时内就会派单处置，复杂问题会有遵义院的技术人员专门协调，不用为维修跑断腿。当然，日常的建筑维修机制的主要目的还是确保我们世界的建筑在100年后依然有安全保障。\\
维修检测机器人（参考：ANYbotics ANYmal）https://www.anybotics.com

\section{10.3 网络与数据安全}\label{ux7f51ux7edcux4e0eux6570ux636eux5b89ux5168}

网络安全核心系统全天运行，保证数据和治理不中断，但数量和功耗受到明确限制。AI负责发现问题、给出建议，关键处置仍由制度和人工决定。这种设计防止安全系统本身变成高能耗、高风险的``新中心''。

\subsection{10.3.1 智能防护架构：特色 AI 自适应防火墙}\label{ux667aux80fdux9632ux62a4ux67b6ux6784ux7279ux8272-ai-ux81eaux9002ux5e94ux9632ux706bux5899}

这是我们网络安全的核心特色 ------AI 自适应防火墙能根据实时数据流智能调整防护策略，就像给网络穿了件 ``智能防弹衣''。网络安全机器人会 24 小时监测数据流动，任何异常访问都会被及时拦截。我们的核心数据，比如大家的个人信息、决策记录，都存储在公有数据库里，由加密 AI 专门守护。想要访问这些数据，必须经过 ``个人 AI + 行政 AI'' 双重验证，还得有人民大会的授权，从根源上杜绝数据泄露。除了防火墙，我们还有网络安全巡检机器人每天自动审计网络状态，漏洞扫描系统会定期排查安全隐患，发现问题立刻修复并生成审计报告，提交人民大会备案，让每一次防护都有据可查。

相关设备:1.AI 自适应防火墙（模型部署：NVIDIA A100）2. 漏洞扫描系统（Qualys Cloud）3．网络安全巡检机器人（Cisco Secure-X 自动审计）\\
来源：1.https://www.nvidia.com2.https://www.cisco.com 3.https://www.qualys.com

\subsection{10.3.2 风险应对措施}\label{ux98ceux9669ux5e94ux5bf9ux63aaux65bd}

为了防止网络瘫痪影响正常生活，我们部署了 ``本地 + 安全洞穴'' 的双备份系统。一旦遇到网络故障，应急 AI30 秒内就能切换到备用网络，启真院会在 24 小时内上报人民大会处置进展，确保治理、民生服务不中断。我们还明确了规则：禁止私下传输涉密数据，违规行为会由 AI 代表提交人民大会协商，经所有 AI 代表统一意见后再约束，既保证数据安全，也体现 ``AI 代议制决策'' 的公平性。 双备份系统（本地 + 安全洞穴）采用 LTO-9 磁带阵列 来源：https://www.fujifilm.com

\section{10.4 国家级与社会级安全}\label{ux56fdux5bb6ux7ea7ux4e0eux793eux4f1aux7ea7ux5b89ux5168}

国家级与社会级安全是日轮之冕得以存续、稳定运转的核心根基，既承载着守护国家主权完整的底线责任，也维系着社会系统有序运行、居民生活安宁的根本利益。它不是孤立的 ``防御工事''，而是贯穿个人安全、建筑安全、网络安全的底层支撑，更是我们在复杂未知环境中实现长期发展的底气所在。

\subsection{10.4.1 边境与主权防护}\label{ux8fb9ux5883ux4e0eux4e3bux6743ux9632ux62a4}

针对无人愿意自愿戍边的情况，我们组建了 ``安防机器人 + AI 巡逻队'' 的核心防护力量。海岸巡逻机器人 24 小时监测岸线，发现非法船只自动驱离；口岸的人员检疫机器人会完成全流程管控，外来人员想要进入，数据得同步给人民大会的 AI 代表，协商一致后才能准入，守住国家主权底线。如果遇到特殊情况（非常时期），现有机器人力量不足时，会启动义务服兵役制度，但服役人员主要负责 AI 设备的辅助操作和监控，不用直接参与一线对抗，既保障边境安全，也照顾大家的意愿。

相关机器人：1.海岸巡逻机器人（参考：EMILY 救援机器人）2.人员检疫机器人（参考：UBTech AIMBOT）\\
来源：1.https://emarinesystems.com/emily/ 2.https://www.ubtrobot.com/en/solutions/

\subsection{10.4.2 社会安全治理}\label{ux793eux4f1aux5b89ux5168ux6cbbux7406}

居民之间有矛盾，不用慌，AI 代表会把问题提交给人民大会协商，行政 AI 会按照统一意见调解，公平公正不偏袒。启真院通过 ``长征'' 中枢监控物资集散点，物流机器人和监管 AI 一起保障物资分配公平，杜绝特权占用，让公有制的分配秩序不被破坏。而对于核能设施、钻井平台这些关键场所，玉湖院会实施 24 小时 AI 值守；大家在 Workshop 使用危险设备时，需要先经启真院审批，避免操作风险，让生产生活都在安全框架内进行。

相关机器人：1.物流监管机器人（参考：Kiva/AMR 自动仓储系统）2.核设施监控机器人（TMI inspection bot 参考设计）\\
来源：1.https://www.amazonrobotics.com 2.https://www.iaea.org/topics/nuclear-safety

\subsection{10.4.3 武装力量及军事装备}\label{ux6b66ux88c5ux529bux91cfux53caux519bux4e8bux88c5ux5907}

遵循宪法第二十八条、三十二条``国家武装力量由人民大会统一领导''的原则，武装力量由玉湖院（国防安全院）统一建设、训练与调度，所有武装装备由人民大会每季度审核，确保符合``防御性武装，不主动扩张''的国家战略定位。相关武装力量有以下几部分组成：

（一）AI主导的无人化防御力量

\begin{enumerate}
\def\labelenumi{\arabic{enumi}.}
\item
  自主海岸防御无人机群（Swarm Defense Drones）\\
  用于海上主权巡逻、发现非法进入、紧急拦截。无人机群与``长征''中枢直连，自动规划巡逻路径，在强风力条件下仍可执行任务，若识别为``高危险来船''，自动广播警告并部署海面拦截浮标。\\
  来源：https://anduril.com
\item
  自动化海面防卫机器人艇（Surface Defense USV）\\
  参考现实技术 Saildrone Voyager + 美国海军USV项目。 在未来世界中的增强能力：具备声学威慑、强光干扰、网状拦阻发射能力；遵循``非杀伤性优先''原则，所有作战动作限制在压制、拦截、隔离范围\\
  来源：https://www.saildrone.com
\end{enumerate}

（二）有人+AI协同的自卫武装力量

在有居民义务参军的情况下，由居民志愿军与义务兵役制度补充组成。\\
AI指挥外骨骼护卫队（Exo-Guard Units）\\
参考设备：Sarcos Guardian XO 外骨骼\\
来源：https://www.sarcos.com\\
给义务服役人员配备外骨骼，提升机动性和负载能力，主要用于建筑、核能场站的警戒，也能在灾后转为救援搬运，民用军用灵活切换。

（三）以人为主导，AI辅助控制的导弹系统

一．地面中程拦截系统（MR-SHIELD 系统）：MR-SHIELD-30 中程截击导弹

参考现实技术：以色列``铁穹''Iron Dome/美制 NASAMS\\
并加入未来版 AI 控制与非致命性策略。具有拦截中小型无人机群、巡航器、武装船只发射的威胁物的防御用途。安全设计：导弹采用``可控破片''设计，只在高空爆破；不落入地面，不对居民区造成二次伤害。 来源：1.https://www.rafael.co.il2.https://www.kongsberg.com

二．导弹武器伦理锁系统（Ethical Lock-M2）

未来科技中，导弹将全部内置伦理锁 M2（Ethical Lock Mechanism 2nd Gen）：\\
1.不可对非来袭方向目标锁定；2.不得锁定民用船只、民航、外来科研船舶；3.不可执行主动侦杀任务；4.攻击指令需三重授权：人类双确认；行政AI认定；人民大会实时记录；5.导弹内置``杀伤禁区限制程序''：凡是落点在人口密集区 → 自动失效；凡是风险不确定 → 自动切换``非爆破拦阻模式''

三．导弹武器与``长征''中枢协同机制

所有导弹系统均与``长征''中枢进行三级数据交换： 1. 战场实时建模（0.1 秒级）：AI 重建三维战场图，预测目标未来位置。 2. 风险等级判断：仅当目标被识别为：攻击性飞行物/未响应警告的高速来船/敌对无人机群，才允许进入第 3 步。 3. 授权链启动：AI 不具备``自动开火''权限，必须经三重授权。

\section{10.5 灾害应急预案}\label{ux707eux5bb3ux5e94ux6025ux9884ux6848}

面对地震、台风、生态灾害这些外界威胁，我们的核心是 ``早预判、快响应、稳恢复''，不搞过度干预，与自然环境保持可控关系。

\subsection{10.5.1 组织与预警体系}\label{ux7ec4ux7ec7ux4e0eux9884ux8b66ux4f53ux7cfb}

``长征'' 中枢是我们的应急指挥部，整合了所有灾害数据，一旦发现风险，会通过 AI 代表和公共终端，1 分钟内把预警信息推送给每一个人。应急小组由 ``玉湖院机器人编队 + 居民技能者'' 组成，启真院统一调度分工，让专业的人做专业的事。\\
灾害数据集成系统（参考：IBM The Weather Company）与广播系统\\
来源：1.https://www.ibm.com2.https://lora-alliance.org

\subsection{10.5.2 应急处置流程}\label{ux5e94ux6025ux5904ux7f6eux6d41ux7a0b}

遇到轻度灾害，避难引导机器人会引导大家前往安全避难点；如果是核隐患这类中度以上灾害，玉湖院 15 分钟内就能完成居民转移，启真院会暂停非必要生产，优先保障应急物资。\\
灾后恢复也有明确流程：启真院的仓储 AI 按 ``按需分配'' 发放应急物资；遵义院的工程机器人优先修复核能、通信这些关键设施；建德院会组织 Workshop 的创作者参与修复设计，确保 1 周内就能恢复正常秩序，让大家尽快回归安稳生活。

相关机器人：1.避难引导机器人（参考：Softbank Whiz）2.工程维修机器人（参考：SARcos Guardian XT）\\
来源：1.https://www.softbankrobotics.com2.https://www.sarcos.com

\subsection{10.5.3 针对性风险防控}\label{ux9488ux5bf9ux6027ux98ceux9669ux9632ux63a7}

针对地质灾害，玉湖院部署了地震监测 AI，建筑都避开断层带，灾害发生后破拆 AI 会快速开辟救援通道，启真院会暂停核能发电切换备用能源，把安全放在第一位。\\
面对台风这类气候灾害，气象 AI 会实时监测路径，遵义院的机器人会提前加固海堤，海上钻井会自动关闭；内涝区域有排水机器人抽排，大家会收到 AI 的转移指令，不用慌慌张张乱躲藏。\\
对于外来入侵生物，宜山院的生态机器人每天都会扫描排查并清除，也禁止大家私自带外来物种进入，守住我们的生态链，让生活环境始终宜人。

相关设备：1.地震监测 AI（参考：中国地震局宽频带监测台站）2.破拆机器人（参考：Brokk 70）3.海堤加固机器人（参考：Built Robotics RPD）4.排水机器人（参考：Sewer Robotics）5.生态监测机器人（参考：AgEagle eBee X）\\
来源：1.https://www.sewerrobotics.com2.https://www.builtrobotics.com3.https://www.brokk.com4.http://www.cea.gov.cn5.https://www.ageagle.com

\section{10.6 未来安全与应急管理：适配扩容，智能进化}\label{ux672aux6765ux5b89ux5168ux4e0eux5e94ux6025ux7ba1ux7406ux9002ux914dux6269ux5bb9ux667aux80fdux8fdbux5316}

随着世界扩容与技术迭代，安全与应急管理将向 ``主动预判、柔性适配、协同共生'' 升级，在守住安全底线的同时，兼顾自由与发展。

``长征'' 中枢将整合长期数据，提前识别健康隐患、建筑疲劳、灾害趋势等潜在风险，在萌芽阶段调度机器人完成干预。AI 与机器人形成 ``感知 - 分析 - 处置'' 闭环，安全机器人动态调整巡检路线，医疗辅助机器人强化诊断与资源对接能力，让风险无需干预即被化解。\\
采用 ``模块化扩容 + 差异化防护'' 模式，新增区域按功能定制安全方案（居住区侧重人身安全、产业区强化设备防护）。

支持安全需求个性化定制，居民可调整监测阈值、机器人响应距离；义务兵役优化为技能适配岗位，避免直面危险，让责任承担更贴合意愿。

联合其他世界搭建协同平台，共享灾害应对经验、应急物资与技术支持，实现跨世界风险共防。网络安全升级为 ``全域联防''，核心数据多世界分布式备份；统一安全标准与应急流程，提升协同响应效率，最终达到提高安全保障的效果。
